\def\type{}
\input{header}
\title{The Langlands--Deligne local constant}
\author{David Kurniadi Angdinata}
\date{Tuesday, 3 February 2026}

\begin{document}

\maketitle

\section{Motivation}

\emph{L-functions} are complex functions that encode information of objects that are typically arithmetic in nature, which are conjecturally meromophic over all of $ \C $ and conjecturally satisfy functional equations. For instance, Riemann/Dedekind $ \zeta $-functions and Artin L-functions encode the behaviour of primes in number fields, Dirichlet/Weber/Hecke L-functions encode the coefficients of modular forms, and Hasse--Weil L-functions encode the reduction data of algebraic varieties. They play the role of half of the modern Langlands philosophy, whose web of research spans almost all of modern pure mathematics, and many crucial algebraic and analytic questions remain unanswered. Most notably, the generalised Riemann hypothesis describes the zeroes of the Dedekind $ \zeta $-function of a number field, while the generalised Birch and Swinnerton-Dyer conjecture describes the arithmetic of an abelian variety in terms of its Hasse--Weil L-function. More specifically, the global Langlands conjectures claim that the L-functions of motivic origin should correspond to the L-functions of suitable automorphic forms, which are realised from Galois and automorphic representations respectively. These conjectures vastly generalise the modularity theorem that Hasse--Weil L-functions of rational elliptic curve are Hecke L-functions of certain weight two modular forms.

While much has been developed in the automorphic picture, a good category of motives, which are meant to be cohomological incarnations of algebraic varieties, remains conjectural. Nevertheless, a motive over a global field $ K $ necessary has local \emph{$ \lambda $-adic realisations} for primes $ \lambda $ of some fixed number field $ F $. This is a suitably compatible system of continuous \emph{$ \lambda $-adic representations} $ \rho_{\lambda, v} : G_v \to \GL(V_{\lambda, v}) $ for each place $ v $ of $ K $, where $ G_v $ is the absolute Galois group of $ K_v $ endowed with the profinite topology and $ V_{\lambda, v} $ is a finite-dimensional vector space over $ F_\lambda $ endowed with the $ \lambda $-adic topology. Instead of working with $ \lambda $-adic representations, whose topologies crucially depend on $ \lambda $, Grothendieck and Deligne showed that their category is equivalent to that of local \emph{Weil--Deligne representations} $ (\rho_v, N_v) $, which are topologically agnostic. This is the data of a \emph{Weil representation} $ \rho_v : W_v \to \GL(V_v) $, which is a continuous representation of the \emph{Weil group} $ W_v \le G_v $ of $ K_v $ acting on a finite-dimensional complex vector space $ V_v $ endowed with the discrete topology, and a certain endomorphism $ N_v : V_v \to V_v $. The Weil group $ W_v $ contains the inertia subgroup $ I_v \le G_v $ and is endowed with the finest topology where $ I_v \le W_v $ is open. Given a suitably compatible system $ \rho := \{\rho_v\}_v $ of Weil representations, its L-function $ L(\rho) $ is then defined to be the product of the local \emph{Euler factor} $ L(\rho_v) $ for each place $ v $ of $ K $, defined as the inverse characteristic polynomial of a choice of geometric Frobenius at $ v $ under the subrepresentation of $ \rho_v $ invariant under $ I_v $ and evaluated at $ 1 $.

For each place $ v $ of $ K $, fix a Weil representation $ \rho_v : W_v \to \GL(V_v) $, a non-trivial additive character $ \psi_v : K_v \to \C^\times $, and an additive Haar measure $ \mu_v $ on $ K_v $, and assume that they satisfy similar compatibility conditions. Langlands and Deligne independently proved that there is a \emph{local constant} $ \epsilon(\rho_v, \psi_v, \mu_v) $, sometimes also called the \emph{$ \epsilon $-factor} or the \emph{root number}, which is uniquely characterised by certain representation-theoretic properties. Under the prevailing assumptions, their product $ \epsilon(\rho) $ converges and is independent of the choices of non-trivial additive characters and additive Haar measures. For each $ s \in \C $, they also proved that the system $ (\rho, s) := \{\rho_v \otimes |\cdot|_v^s\}_v $ is related to its dual system $ (\rho^\vee, 1 - s) $ by a functional equation
$$ L(\rho, s) = \epsilon(\rho, s) \cdot L(\rho^\vee, 1 - s). $$
When $ \rho $ is one-dimensional, which corresponds to a quasi-character by local class field theory, this result was proven by Hecke and modernised into the framework of harmonic analysis in Tate's thesis.

The aim of this study group is to understand the existence statement of the Langlands--Deligne local constant $ \epsilon(\rho_v, \psi_v, \mu_v) $ when $ v $ is a non-archimedean place of $ K $ and sketch a proof if time permits.

\pagebreak

\section{Talks}

Here are some topics to be covered in the study group.
\begin{itemize}
\item Local fields
\begin{itemize}
\item Basic definitions: absolute values, discrete valuations, local compactness, normalised uniformisers, valuation rings, residue fields [Ser79, Chapter I Section 1] [CF80, Chapter I Section 1]
\item Ramification theory: Frobenius substitutions, inertia subgroups, Weil groups, higher ramification groups [Ser79, Chapter I Section 8] [CF80, Chapter I Section 5 to Section 9]
\item Fourier analysis: additive characters, Haar measures, Schwarz--Bruhat functions, Fourier transforms [CF80, Chapter XV Section 2.2 and Section 2.3] [Poo15, Section 4.2 to Section 4.5]
\end{itemize}
\item Local Weil representations
\begin{itemize}
\item Weil groups in local class field theory [Tat97, Section 1.1 to Section 1.3]
\item Quasi-characters and representations of Galois type [Tat97, Section 2.2]
\item Inductive functions of representations [Tat97, Section 2.1 and Section 2.3]
\item L-functions and conductors [Roh93, Section 8 and Section 10] [Tat97, Section 3.1 and Section 3.3]
\end{itemize}
\item Local constants
\begin{itemize}
\item Existence, uniqueness, properties [Del75] [Roh93, Section 11] [Tat97, Section 3.2 and Section 3.4]
\item Abelian local constants [CF80, Chapter XV Section 2.4] [RV99, Section 7.1] [Poo15, Section 4.9]
\item Non-abelian local constants [Lan70] [Del73, Section 4] [Tat77, Section 2] [BH06, Section 30]
\end{itemize}
\end{itemize}
The main reference will be the sections up to Theorem 3.4.1 of Tate's 1997 article [Tat97].

\section{References}

The existence of the Langlands--Deligne local constant is surveyed in various references.
\begin{itemize}
\item[Del75] Deligne's 1975 article \emph{Les constantes des \'equations fonctionnelles}
\item[Roh93] Rohrlich's 1993 article \emph{Elliptic curves and the Weil--Deligne group}
\item[Tat97] Tate's 1997 article \emph{Number theoretic background}
\end{itemize}
The proofs of existence by Langlands and Deligne are detailed in various references.
\begin{itemize}
\item[Lan70] Langlands's 1970 lecture notes \emph{On the functional equation of the Artin L-functions}
\item[Del73] Deligne's 1973 article \emph{Les constantes des \'equations fonctionnelles des fonctions L}
\item[Tat77] Tate's 1977 article \emph{Local constants}
\item[BH06] Bushnell and Henniart's 2006 book \emph{The local Langlands conjecture for GL(2)}
\end{itemize}
Here are some references for basic algebraic number theory up to and including Tate's thesis.
\begin{itemize}
\item[Wei67] Weil's 1967 book \emph{Basic number theory}
\item[Ser79] Serre's 1979 book \emph{Local fields}
\item[CF80] Cassels and Fr\"ohlich's 1980 book \emph{Algebraic number theory}
\item[RV99] Ramakrishnan and Valenza's 1999 book \emph{Fourier analysis on number fields}
\item[Poo15] Poonen's 2015 lecture notes \emph{Tate's thesis}
\end{itemize}

\end{document}